\documentclass[a4paper,11pt]{article}
\usepackage[utf8]{inputenc}
\usepackage[english]{babel}
\usepackage[margin=2.5cm]{geometry}
\usepackage{amsmath}
\usepackage{amssymb}
\usepackage{graphicx}
\usepackage{xcolor}
\usepackage{enumitem}
\usepackage{booktabs}

\title{\textbf{TAF ISC - UE AIST\\Examination, December 2023\\part B (1h30)}}
\date{}

\begin{document}

\maketitle

\textbf{NAME:} \rule{8cm}{0.4pt}

\section*{QUESTIONS}

\begin{enumerate}
    \item Cite by order of importance at least three advantages of the silica optical fiber with respect to other wireline transmission media.
    
    \vspace{2cm}
    
    \item What is the line bit rate of a 43.7 Gbaud Polarization Multiplexed 64-QAM transponder?
    
    \vspace{2cm}
    
    \item What is the Nyquist spectral occupancy of the signal transmitted by such a transponder?
    
    \vspace{2cm}
    
    \item What is the distance limiting factor for very long optically amplified transmission systems?
    
    \vspace{2cm}
    
    \item Which technology is used at the digital level to improve the transmission performance (reach and/or bit error rate) of modern optical transmission systems?
    
    \vspace{2cm}
    
    \item Which typical link budget (in dBs) do you need to transmit over 100 km of silica fiber?
    
    \vspace{2cm}
    
    \item Why does coherent detection improve capacity with respect to previous IM/DD systems?
    
    \vspace{2cm}
    
    \item Do you still need in-line dispersion compensating devices in optical coherent systems?
    
    \vspace{2cm}
\end{enumerate}

\newpage

\section*{PROBLEM}

An 8260 km long submarine cable system has to be deployed. The lineic fiber attenuation is 0.18 dB per km in the C band (4 THz wide). The regularly spaced EDFA optical amplifiers have a noise figure of 5 dB, a maximum output power of 14.3 dBm and adjustable gain up to 12.6 dB.

A multiplex of 100Gbit/s DP-8QAM channels is transmitted through the line. The channel spectral occupancy is 25 GHz and the target OSNR at the receiver end is 15 dB.

\textbf{NB:} The optical reference band is 12.5 GHz and OSNR is considered as infinite at the transmitters' output.

In order to minimize cost, the strictly sufficient number of amplifiers will be installed.

\subsection*{Questions:}

\begin{enumerate}
    \item How many amplifiers are installed on the line and what is the amplifier spacing?
    
    \vspace{3cm}
    
    \item What is the maximum number of channels that can be transmitted?
    
    \vspace{3cm}
    
    \item Which channel spacing do you choose?
    
    \vspace{3cm}
    
    \item What is the channel spectral efficiency?
    
    \vspace{3cm}
    
    \item What is the system spectral efficiency?
    
    \vspace{3cm}
    
    \item How could we increase the system capacity without modifying the installed subsea line?
    
    \vspace{3cm}
    
    \item What is the maximum system capacity that can be achieved using the installed subsea cable when transmitting 37.5 GHz wide DP-QPSK 100 Gbit/s channels with 13 dB OSNR target?
    
    \vspace{3cm}
\end{enumerate}

\end{document}
